
\subsection{PSICROMETRÍA}

De la ecuación del estado que describe el comportamiento de gases:
\eqEosWithZ

Para una mezcla gaseosa (componente  vapor y componente gas), se puede escribir las ecuaciones

\eqDoubleEosVaporGas

Dividiendo la ecuación \eqref{eq:eosVapor} y la ecuación \eqref{eq:eosGas}:
\begin{align}
    \frac{{\pVapor}{\volume[]{}[]}}{{\pGas}{\volume[]{}[]}} &= \frac{{\zVapor}{\vaporMole}{\universalConstantGas}{\tVapor}}{{\zGas}{\gasMole}{\universalConstantGas}{\tGas}}
\end{align}

Simplificando:

\begin{align}
    % Simplificamos el Volumen con el índice [1] (Rojo)
    % Simplificamos la constante R con el índice [2] (Azul)
    \frac{P_v \cancelC[1]{V}}{P_g \cancelC[1]{V}} &= \frac{z_v n_v \cancelC[2]{R} T_v}{z_g n_g \cancelC[2]{R} T_g}
\end{align}

Queda,

\begin{align}
    \frac{\pVapor}{\pGas} &= \frac{\zVapor \vaporMole \tVapor}{\zGas \gasMole \tGas}
\end{align}

Según la ley de Willard Gibbs, las variables independientes del sistema, son: \pTotal , \temperature[]{} y \pVapor (o \vaporMoleFraction)

\eqPhaseRule

De acuerdo a la Ley de Dalton,

\eqTotalPressure


% \subsubsection{Temperatura de punto de rocío (\tDewPoint)}
% \subsubsection{Temperatura de punto de rocío (\texorpdfstring{\tDewPoint}{DewPoint}) }
\subsubsection{Temperatura de punto de rocío (\texorpdfstring{\tDewPoint}{Troc})}

Sabemos que temperatura de rocío es la temperatura a la que la presión de vapor (\pVapor), iguala la presión de saturación; es decir;
\begin{align}
    \pVapor = \pSat(\tDewPoint) \quad \rightarrow \quad \tDewPoint = \pSatInv(\pSat)
\end{align}

Donde $\pVapor$ se puede calcular a partir de la ecuación de Antoine, Wagner, Lee y Kesler; sin embargo, se usará la biblioteca de \href{https://coolprop.org}{\textcolor{purple}{CoolProp}}, el cual es utilizado por Thermocycle, \href{https://pdsim.sourceforge.net/}{\textcolor{purple}{PDSim}}, \href{https://achp.sourceforge.net/ACHPGUI/Tutorial.html#getting-started}{\textcolor{purple}{ACHP}}, \href{https://sourceforge.net/projects/dwsim/}{\textcolor{purple}{DWSim}}, StateCalc, SmoWeb, T-Props, \href{https://app.fproperties.org/}{\textcolor{purple}{fProperties}}, \href{https://github.com/dvd101x/web-thermodynamics}{\textcolor{purple}{CoolPropJavascriptDemo}}, \href{https://app.psolver.org/app/}{\textcolor{purple}{pSolver}}.

\subsubsection{Humedad relativa (\texorpdfstring{\relativeHumidity[]{}}{phi})}
    \eqnRelativeHumidity



\subsubsection{Relación de humedad o humedad absoluta}
    \eqAbsoluteHumidity
    Expresándolo en relación de presiones del vapor y el gas:
    \eqMassToMolOfVaporAndGas
    
    Reemplazando las ecuaciones \eqref{eqn:mass_to_mole_vapor} y \eqref{eqn:mass_to_mole_gas} en la ecuación \eqref{eq:AbsoluteHumidity}
    \eqAbsoluteHumidityToMol
    
    Del número total de moles,
    \eqTotalMoles*

    Tenemos la suma de las fracciones molares    
    \eqGasFracAsVaporFrac

    \eqAbsoluteHumidityToMolAsMolarFraction

    \eqAbsoluteHumidityToMolAsMolarVaporFraction


    La fracción molar del vapor puede escribirse asi,
    \eqMoleVaporFractionAsAbsoluteHumidity

    Reemplazando \eqref{eq:MoleVaporFractionAsAbsoluteHumidity} en la ecuación \eqref{eq:GasFracAsFrac} tenemos la fracción molar para el gas como una función de la humedad absoluta,
    \eqMoleGasFractionAsAbsoluteHumidity

    También se puede escribir en función de la presión de vapor,
    \eqMixSimplAbsoluteHumidityAsPtotal

    O bien, como una función de la presión de vapor saturado,
    \eqMixSimplAbsoluteHumidityAsPtotalAndPressureSaturated


    El peso molecular promedio se puede calcular a partir de la expresión,
    \eqMolecularWeightAverage

    
\subsubsection{Volumen específico y molar del aire}
    De la ecuación \eqref{eq:EosWithZ}; sustituimos 
    \eqEosWithMolarVolumeZ*
    
    Recordemos que $\molarVolume = \V / \n$; despejamos $\p$,
    \eqEosPressureWithMolarVolumeZ*
    
    % Y $\z$ para la ecuación virial, se puede escribir como:
    % \eqZVirialThirdCoeff*
    % \eqVirialSecondMix
    % \eqVirialThirdMix
    Para los gases ideales, el factor de compresibilidad toma el valor de $\z = 1$.

    \eqEnthalpyRigorous

    De la ecuación \eqref{eq:enthalpy:rigorous}, considerando que $\zVapor = \zGas = 1$ y  $\molarVolume \to \infty$; se tiene una aproximación para gases ideales, donde los términos viriales se anulan y la entalpía molar depende únicamente de la temperatura y la composición; tal como sigue,
    
    \eqEnthalpyIdealMolar
    
    Además,
    \eqEnthalpyConversion
    
    Las fracción molar $\vaporMoleFraction$ se calcula a partir de la ecuación \eqref{eq:MoleVaporFractionAsAbsoluteHumidity} y $\gasMoleFraction$ desde la ecuación \eqref{eq:MoleGasFractionAsAbsoluteHumidity}; finalmente se tiene la expresión,
    
    
    \eqEnthalpyExpanded*
    Reordenando,
    \eqEnthalpyExpandedDevelop*
    Haciendo un cambio de variable para,
    \eqCpHumid*
    Finalmente tenemos,
    \eqEnthalpyCompact*
    
\subsubsection{Humedad absoluta molar}
    Se define la humedad absoluta molar como:
    \eqnMolarAbsoluteHumidity 